\section{The Inherited Vocabulary: Isidore's Book on Laws}

When the medieval canonists wanted definitions, they did not go first to the Roman jurists directly, and not yet to Aristotle. They went to a bishop of Seville who died in 636. Book V of Isidore's \emph{Etymologies}, \latin{De legibus et temporibus}, is a compressed digest of late Roman jurisprudence---Isidore is transmitting, not inventing---and its opening chapters supplied the axioms from which both Gratian's \emph{Decretum} and Aquinas's treatise on law visibly work.\endnote{Isidore of Seville, \emph{Etymologiae} V, quoted throughout this section from W.~M. Lindsay's critical edition (\emph{Isidori Hispalensis episcopi Etymologiarum sive Originum libri XX}, Oxford: Clarendon Press, 1911), vol.~1, in the scan of the University of Toronto exemplar (Internet Archive item \texttt{isidori01isiduoft}). The three page images carrying the quoted chapters (leaves 192, 193, 196 of the scan) were fetched through the archive's IIIF service, hashed, retained, and read visually; the archive's OCR text served as a discovery aid only. Chapter numbers follow Lindsay. Isidore's own sources here are Roman juristic and rhetorical commonplaces (Gaius, Ulpian, Cicero, the \emph{Institutes} tradition), received at second hand; this article cites Isidore as the transmitter the canonical tradition actually read, not as an authority on classical Roman law.} Four of those chapters carry almost the whole vocabulary this article traces.

\subsection{Divine and human, \latin{fas} and \latin{ius}}

The second chapter of Book V is a single breath: \latin{Omnes autem leges aut divinae sunt, aut humanae. Divinae natura, humanae moribus constant; ideoque haec discrepant, quoniam aliae aliis gentibus placent}---all laws are either divine or human; divine laws stand by nature, human laws by usages, and that is why human laws diverge, since different things please different peoples. Then a pair of technical terms: \latin{Fas lex divina est, ius lex humana. Transire per alienum fas est, ius non est}---\latin{fas} is divine law, \latin{ius} human law; to cross another's land is \latin{fas}, it is not \latin{ius}.\endnote{Etym.\ V.ii, read at the retained image of leaf 192. The example is exact and illuminating: no divine ordinance forbids walking across a field; positive human law, protecting property, does. Note what Isidore's shorthand does \emph{not} yet distinguish: ``divine'' law standing \latin{natura} is what the scholastics will call \emph{natural} law, and Isidore has no separate slot here for divinely \emph{revealed} positive law. Gratian will inherit exactly this compression (section 7).}

Everything the later tradition will complicate is already here in embryo, including the imprecision that will have to be repaired. Isidore's ``divine'' law is what stands \latin{natura}---by nature. The scheme has two members, and the divine member is simply the natural. Where in this scheme is the law God gave by speaking---the Decalogue, the Gospel precepts, the sacramental economy---which does not stand by nature and which no human legislator enacted? Isidore does not say; the first sentence of the book, naming Moses as the first to set forth \latin{divinas leges} in writing, quietly assumes it. A taxonomy with a missing drawer is an invitation to successors, and both Gratian and Thomas accepted it in different ways.

\subsection{\latin{Ius}, \latin{lex}, \latin{mos}: written law and custom}

Chapter three sorts the words themselves. \latin{Ius} is the general name; \latin{lex} is a species of it: \latin{Ius generale nomen est, lex autem iuris est species}. And then the definition the Middle Ages memorized: \latin{Lex est constitutio scripta}---law is a written enactment. \latin{Mos} is \latin{vetustate probata consuetudo, sive lex non scripta}---custom proved by long age, unwritten law; for \latin{lex} takes its name from \latin{legere}, to read, because it is written. Custom (\latin{consuetudo}) is \latin{ius quoddam moribus institutum, quod pro lege suscipitur, cum deficit lex}---a kind of law established by usages, which is taken for law where written law fails; nor does it matter whether it rests on writing or on reason, since reason commends law too.\endnote{Etym.\ V.iii, read at the retained image of leaf 192; the parallel text at Etym.\ II.x carries the same definitions in the rhetoric book. The final clause continues: if law stands by reason, then law will be whatever stands by reason, \latin{dumtaxat quod religioni congruat, quod disciplinae conveniat, quod saluti proficiat}---provided it agrees with religion, befits discipline, profits salvation.}

Three commitments in that chapter shaped a millennium of canon law. Written enactment is the paradigm of positive law but not its whole extent: custom is genuinely law, not tolerated fact. Custom's title to be law is subsidiary---\latin{cum deficit lex}, where enacted law fails. And the deepest criterion is neither writing nor age but reason. Canon 24~§2 of the present Code, requiring that a custom be \latin{rationabilis} to obtain force of law, is this sentence of Isidore still working, fourteen centuries on (section 9).

\subsection{Natural, civil, and the law of nations}

Chapter four opens the tripartite division the Roman jurists had made famous: \latin{Ius autem naturale [est], aut civile, aut gentium}---law is natural, or civil, or of nations. Natural law is \latin{commune omnium nationum}, held everywhere \latin{instinctu naturae, non constitutione aliqua}---by the prompting of nature, not by any enactment: the union of man and woman, the succession and rearing of children, the common possession of all things and the one liberty of all, the acquisition of what is taken from sky, land, and sea; the restitution of a thing deposited or money entrusted, the repelling of violence by force. \latin{Nam hoc, aut si quid huic simile est, numquam iniustum, sed naturale aequumque habetur}---for this, and whatever resembles it, is never unjust, but is held to be natural and equitable.\endnote{Etym.\ V.iv, heading and opening words at the foot of leaf 192, body read at the retained image of leaf 193. The list is Roman (compare Ulpian and Gaius in the \emph{Digest}'s opening titles, which this article has not independently examined and cites only as the well-known background); its tensions are Roman too---``common possession of all things'' and ``one liberty of all'' sat in the same legal culture that enforced property and slavery, a tension the scholastics resolved by distinguishing what nature \emph{institutes} from what it merely does not forbid, and which Gratian's D.8 exploits directly (section 7).} Civil law is what each people or city establishes as its own, \latin{humana divinaque causa}; the law of nations---occupations of territory, war, captivity, treaties, embassies---is so called \latin{quia eo iure omnes fere gentes utuntur}, because nearly all nations use it.\endnote{Etym.\ V.v--vi, read at the retained image of leaf 193. Note that Isidore's \latin{ius gentium} is a category of \emph{human} practice common to peoples, not modern international law; nothing in this article depends on conflating them.}

The important thing for what follows is the fault line inside the scheme. ``Natural'' names a source (nature, not enactment); ``civil'' names an enactor (this people); ``of nations'' names an extent (nearly everyone). These are not parallel criteria, and the tradition had to decide which cut was fundamental. Aquinas's answer---that the deep division is between what is derived from natural law as \emph{conclusion} and what is derived as \emph{determination}, with the \latin{ius gentium} on one side of that line and civil law on the other---will occupy section 5.

\subsection{What a law must be}

Chapter twenty-one is a checklist, and it became the most quoted sentence of the book: \latin{Erit autem lex honesta, iusta, possibilis, secundum naturam, secundum consuetudinem patriae, loco temporique conveniens, necessaria, utilis, manifesta quoque, ne aliquid per obscuritatem in captionem contineat, nullo privato commodo, sed pro communi civium utilitate conscripta}---law will be honorable, just, possible, according to nature, according to the custom of the country, suited to place and time, necessary, useful, manifest too---lest through obscurity it contain a trap---written for no private advantage, but for the common benefit of the citizens.\endnote{Etym.\ V.xxi, read at the retained image of leaf 196. The immediately preceding chapter (V.xx, on the purpose of laws: that human audacity be checked by fear of them, that innocence be safe among the wicked) reaches this article only as Gratian's D.4 pr.\ quotes it (section 7); it was not separately verified at the page image, and nothing turns on its wording.}

Every phrase of that sentence had a career. \latin{Possibilis, secundum naturam, secundum consuetudinem patriae} becomes Aquinas's argument that human law must not try to forbid every vice (I-II q.~96 a.~2, quoting this very chapter). \latin{Pro communi utilitate civium conscripta} becomes the ``common good'' clause of the scholastic definition of law. And the whole checklist, received through Gratian's D.4, is the ancestor of the modern canonical instinct that a law failing these tests is a candidate for non-reception, desuetude, or emendation rather than reverence. Isidore wrote an encyclopedia entry; he turned out to have written the specification sheet against which the Latin Church still audits its own legislation.

One boundary should be stated before leaving him. Isidore carries no theory. There is no account here of \emph{why} nature obliges, of how divine speech relates to natural instinct, of what happens when the categories conflict. The \emph{Etymologies} is a filing cabinet, and its authority in what follows is the authority of a filing cabinet everyone used: it fixed the labels. The theory came six centuries later, and it is to that we now turn.
